% !TEX root = ../main.tex
% Section 6. Page budget: part of the 4.25 pages for §5-§10 (report/README.md).
% Filled by #75. Stage 1 (this version): the network baseline only, from data/baseline/ (#70,
% #83). Stage 2 adds the §11 experiments from #77 (analysis/out/, #71).
\section{Performance measurement, results and graphs (R6)}
\label{sec:performance}

\subsection{Network baseline}
\label{sec:baseline}

Before measuring meshsync we measured the network itself on the demo LAN (14 November 2026), in
each direction between the three hosts: round-trip time with \code{ping} (5 runs of 10 echoes;
each run's average is one sample), meshsync's own PING/PONG round trips from a 3-minute
plain-mode run (Table~\ref{tab:rtt}), and the TCP throughput ceiling with \code{iperf3} (5 runs of
10\,s, receiver side, 1\,Mbit = $10^6$ bit; Table~\ref{tab:iperf}). Values are mean $\pm$ sample
standard deviation; ``src $\rightarrow$ dst'' is the side that sends the echo request or the data.
% data/baseline/README.md ("Files", method settings, "Results"); summary.csv; pingpong.csv.

\begin{reptable}{>{\raggedright\arraybackslash}p{\dimexpr0.30\linewidth-2\tabcolsep\relax}>{\raggedright\arraybackslash}p{\dimexpr0.35\linewidth-2\tabcolsep\relax}>{\raggedright\arraybackslash}p{\dimexpr0.35\linewidth-2\tabcolsep\relax}}
\caption{Round-trip time between the hosts (ms).}\label{tab:rtt}\\
\toprule
\colhead{src $\rightarrow$ dst} & \colhead{\code{ping} (n = 5 runs)} & \colhead{meshsync PING/PONG (n)} \\
\midrule
\endfirsthead
\multicolumn{3}{@{}l}{\tablecontinued}\\
\toprule
\colhead{src $\rightarrow$ dst} & \colhead{\code{ping} (n = 5 runs)} & \colhead{meshsync PING/PONG (n)} \\
\midrule
\endhead
\bottomrule
\endlastfoot
kratos $\rightarrow$ defalt & 9.70 $\pm$ 10.27 & 10.91 $\pm$ 5.59 (20) \\
defalt $\rightarrow$ kratos & 8.55 $\pm$ 3.42 & 10.87 $\pm$ 4.81 (20) \\
kratos $\rightarrow$ minipc & 68.02 $\pm$ 25.97 & 71.82 $\pm$ 29.38 (21) \\
minipc $\rightarrow$ kratos & 14.24 $\pm$ 10.78 & 13.19 $\pm$ 3.56 (21) \\
minipc $\rightarrow$ defalt & 17.34 $\pm$ 13.81 & 17.15 $\pm$ 14.62 (21) \\
defalt $\rightarrow$ minipc & 56.53 $\pm$ 5.61 & 35.98 $\pm$ 20.42 (20) \\
\end{reptable}
% ping: data/baseline/summary.csv (metric rtt_ms, n 5). PING/PONG: data/baseline/pingpong.csv,
% mean and sample sd per (host, peer) computed from its rtt_ms rows (n 20-21).

\begin{reptable}{>{\raggedright\arraybackslash}p{\dimexpr0.4\linewidth-2\tabcolsep\relax}>{\raggedright\arraybackslash}p{\dimexpr0.6\linewidth-2\tabcolsep\relax}}
\caption{TCP throughput ceiling, \code{iperf3} (n = 5 runs).}\label{tab:iperf}\\
\toprule
\colhead{src $\rightarrow$ dst} & \colhead{Throughput (Mbit/s)} \\
\midrule
\endfirsthead
\multicolumn{2}{@{}l}{\tablecontinued}\\
\toprule
\colhead{src $\rightarrow$ dst} & \colhead{Throughput (Mbit/s)} \\
\midrule
\endhead
\bottomrule
\endlastfoot
kratos $\rightarrow$ defalt & 16.58 $\pm$ 0.78 \\
defalt $\rightarrow$ kratos & 7.66 $\pm$ 1.26 \\
kratos $\rightarrow$ minipc & 14.29 $\pm$ 3.62 \\
minipc $\rightarrow$ kratos & 27.59 $\pm$ 1.41 \\
minipc $\rightarrow$ defalt & 10.45 $\pm$ 2.63 \\
defalt $\rightarrow$ minipc & 10.38 $\pm$ 1.34 \\
\end{reptable}
% data/baseline/summary.csv (metric throughput_mbps, n 5).

\textbf{Interpreting the RTT.} Each round trip is the sum of the nodal delays
$d_\mathrm{proc} + d_\mathrm{queue} + d_\mathrm{trans} + d_\mathrm{prop}$ on every hop, and here every
packet crosses the air twice (sender to access point to receiver) on one shared channel. On one
LAN $d_\mathrm{prop}$ is negligible, and $d_\mathrm{trans}$ is small for short echo frames, so the
spread comes from queueing and waiting for the channel. The data show this: the fastest single
echo in any run took 2.1\,ms, run averages reach 98.6\,ms, single echoes up to 348.8\,ms, and the
standard deviations are often comparable to the means. Three runs lost one of their ten echoes. The
clearest pattern is that traffic \emph{to} minipc is slow (57--68\,ms) while traffic \emph{from}
it is fast (14--17\,ms): Wi-Fi power saving is on for minipc, so the access point holds frames
for it until its radio wakes. That waiting is queueing delay, not distance. meshsync's PING/PONG
agrees with \code{ping} within 4\,ms in five of six directions, usually slightly above it,
which fits our own extra $d_\mathrm{proc}$ (Python, framing, the reader thread). The exception, defalt $\rightarrow$
minipc (36 vs 57\,ms), also waits on minipc's sleep cycle, and the high standard deviation (20\,ms)
reflects when each probe happened to arrive. The \code{iperf3} ceiling is 7.7--27.6\,Mbit/s and
depends on direction (defalt sends at 7.7 but receives at 16.6\,Mbit/s): this is the upper bound
for any meshsync transfer between these hosts.
% data/baseline/rtt.csv: min(min_ms) 2.077 (defalt->kratos run 2), max(avg_ms) 98.608
% (kratos->minipc run 5), max(max_ms) 348.817 (minipc->defalt run 4), loss_pct 10 in 3 runs.
% Power save and the defalt->minipc explanation: data/baseline/README.md ("Anomalies").
% PROJECT_CONTEXT.md §11 (delay model).

\subsection{Experiments}
\label{sec:experiments}

\textbf{Method.} We ran the §11 conditions on the same three hosts (18--19 November 2026), with
kratos as the driver and source and minipc and defalt as receivers, each condition 5 times. The
workload is set A: a 100\,MiB file (104\,857\,600 bytes), then a 1\,KiB overwrite (delta) or a 1\,KiB
insert at offset 0 (fixed vs content-defined chunking). Impairment is \code{tc netem} on kratos's
Wi-Fi interface: 5\,\% loss, or 50\,ms added delay. For each run, the receivers' values are averaged
first; we then report mean $\pm$ sample standard deviation over the n = 5 runs, as
\code{analysis/analyze.py} computes it. Throughput is chunk bytes divided by transfer time, in MB/s
($10^6$ bytes). One TLS row that \code{analyze.py} flags as a duplicate fetch (\#60) is excluded.
Table~\ref{tab:experiments} gives the results; Figures~\ref{fig:exp-full-delta}
and~\ref{fig:exp-impaired} show two of them. The chunk-size condition (16/64/256\,KiB) was not
run.
% data/README.md (§11 multi-host results: setup, interruptions, tls/run4 flagged);
% scripts/measure/run.py impaired_profiles, LOSS_PCT_DEFAULT 5, DELAY_MS_DEFAULT 50 (hosts.py
% --loss-pct/--delay-ms defaults; data/README.md records no other value); analysis/README.md
% (statistics, throughput). Tables: analysis/out/<experiment>.md from analysis/analyze.py.

\begin{reptable}{>{\raggedright\arraybackslash}p{\dimexpr0.22\linewidth-2\tabcolsep\relax}>{\raggedright\arraybackslash}p{\dimexpr0.12\linewidth-2\tabcolsep\relax}>{\raggedright\arraybackslash}p{\dimexpr0.22\linewidth-2\tabcolsep\relax}>{\raggedright\arraybackslash}p{\dimexpr0.24\linewidth-2\tabcolsep\relax}>{\raggedright\arraybackslash}p{\dimexpr0.20\linewidth-2\tabcolsep\relax}}
\caption{§11 conditions, mean $\pm$ sd over n = 5 runs.}\label{tab:experiments}\\
\toprule
\colhead{Condition} & \colhead{Variant} & \colhead{Transfer time (s)} & \colhead{Chunk bytes received} & \colhead{Throughput (MB/s)} \\
\midrule
\endfirsthead
\multicolumn{5}{@{}l}{\tablecontinued}\\
\toprule
\colhead{Condition} & \colhead{Variant} & \colhead{Transfer time (s)} & \colhead{Chunk bytes received} & \colhead{Throughput (MB/s)} \\
\midrule
\endhead
\bottomrule
\endlastfoot
1 Full vs delta & full & 227.5 $\pm$ 74.99 & 104\,874\,436 & 0.5116 $\pm$ 0.1564 \\
 & delta & 5.204 $\pm$ 10.92 & 75\,667 & 0.3633 $\pm$ 0.3376 \\
2 Fixed vs CDC (insert) & fixed & 143.1 $\pm$ 49.58 & 104\,877\,836 & 0.8572 $\pm$ 0.4003 \\
 & CDC & 0.3006 $\pm$ 0.4061 & 66\,594 & 0.5789 $\pm$ 0.3944 \\
3 Sources & 1 source & 73.33 $\pm$ 30.9 & 104\,894\,652 $\pm$ 45\,205 & 1.66 $\pm$ 0.7107 \\
 & 2 sources & 66.05 $\pm$ 10.22 & 104\,874\,436 & 1.617 $\pm$ 0.2347 \\
5 Link & clean & 116.4 $\pm$ 33.57 & 104\,874\,436 & 1.031 $\pm$ 0.2897 \\
 & 5\,\% loss & 474.2 $\pm$ 143.7 & 104\,874\,436 & 0.2352 $\pm$ 0.05852 \\
 & +50\,ms delay & 118.8 $\pm$ 27.07 & 104\,874\,436 & 0.946 $\pm$ 0.2376 \\
6 Security & plain & 102.9 $\pm$ 16.81 & 104\,874\,436 & 1.055 $\pm$ 0.1655 \\
 & TLS & 90.76 $\pm$ 19.55 & 105\,160\,119 $\pm$ 620\,599 & 1.23 $\pm$ 0.2855 \\
\end{reptable}
% analysis/out/{full_delta,fixed_cdc,sources,impaired,tls}.md (chunk bytes sd 0 except where
% shown; 1 source: 104,894,652 +- 45,205). Connection setup time (tls.md): plain 0.6517 +- 1.078 s,
% TLS 0.2537 +- 0.05434 s.

\begin{figure}[!htbp]\centering
  \begin{minipage}[t]{0.48\linewidth}\centering
    \fitgraphic[0.27\textheight]{exp-full-delta.pdf}
    \caption{Chunk bytes received, full copy vs delta (log scale, n = 5, $\pm$1 sd).}
    \label{fig:exp-full-delta}
  \end{minipage}\hfill
  \begin{minipage}[t]{0.48\linewidth}\centering
    \fitgraphic[0.27\textheight]{exp-impaired.pdf}
    \caption{Transfer time of the 100\,MiB file: clean, 5\,\% loss, +50\,ms delay (n = 5, $\pm$1 sd).}
    \label{fig:exp-impaired}
  \end{minipage}
\end{figure}
% report/figures/exp-full-delta.pdf, exp-impaired.pdf: copies of analysis/out/full_delta.pdf,
% impaired.pdf (analysis/analyze.py).

\textbf{RTT, throughput and the representative operation.} RTT is in Table~\ref{tab:rtt}
(5 runs of \code{ping}, and 20--21 meshsync PING/PONG samples per direction). The representative
operation, syncing the 100\,MiB file to two receivers, took $116.4 \pm 33.6$\,s on the clean link
(condition 5), about $1.03$\,MB/s or 8.2\,Mbit/s. That is about 50--58\,\% of the \code{iperf3}
ceiling from kratos (14.29\,Mbit/s to minipc, 16.58 to defalt; Table~\ref{tab:iperf}). The ceiling is
per pair, and both receivers share one Wi-Fi channel with the source, so it is not additive.

\textbf{Explaining the results with $d_\mathrm{nodal}$.} For the full 100\,MiB transfers,
$d_\mathrm{trans}$ dominates: 104.9\,MB at about 1\,MB/s takes about 100\,s, while a round trip is tens
of milliseconds. The Wi-Fi channel is the bottleneck, and queueing behind it ($d_\mathrm{queue}$)
explains the large run-to-run spread. For delta sync and CDC after an insert, only 75\,667 and 66\,594
chunk bytes cross the network, so $d_\mathrm{trans}$ almost disappears (every run but one under 1.4\,s) and the
remaining time is a few request/reply round trips plus our own $d_\mathrm{proc}$ (hashing, index
compare). 5\,\% loss makes the transfer 4.1 times slower in this batch (2.1 times in a later batch on a better night, Section~\ref{sec:ftp}): lost segments must be retransmitted and TCP
shrinks its congestion window, which shows up as extra waiting, not as distance; the same loss is
behind the 402 retransmissions in Section~\ref{sec:failure}. Adding 50\,ms of delay, a
$d_\mathrm{prop}$-like term, has no clear effect ($118.8 \pm 27.1$ vs $116.4 \pm 33.6$\,s): at
1.031\,MB/s, 50\,ms holds about 52\,KB in flight, far below what our pipelining keeps outstanding
(8 outstanding CHUNK\_REQs per peer $\times$ 64\,KiB average chunk, about 524\,KB), so the link stays busy. Two
sources instead of one give no clear gain ($66.1 \pm 10.2$ vs $73.3 \pm 30.9$\,s): the sources share
the same air time, so the bottleneck does not move. TLS shows no measurable cost: its throughput
($1.23 \pm 0.29$\,MB/s) and setup time ($0.25 \pm 0.05$\,s) are within the spread of plain mode
($1.06 \pm 0.17$\,MB/s, $0.65 \pm 1.08$\,s), because encryption adds a little $d_\mathrm{proc}$ on a link
whose $d_\mathrm{trans}$ is far larger.
% 104874436 B / 1.031e6 B/s = 102 s; 474.2 / 116.4 = 4.07; 1.031e6 x 0.05 = 51,550 B;
% MAX_OUTSTANDING_PER_PEER 8 x CDC_AVG 64 KiB = 524,288 B (constants.py); 1.031 MB/s x 8 = 8.25
% Mbit/s; 8.25/14.29 = 0.58, 8.25/16.58 = 0.50. Per-run values: data/<experiment>/run*.csv.

\textbf{Anomalies.} (1) The same 100\,MiB operation took $227.5 \pm 75.0$\,s in condition 1,
$143.1 \pm 49.6$\,s in condition 2 and $102.9 \pm 16.8$\,s in condition 6: a factor of two between blocks
of one batch. defalt's Wi-Fi was lossy during the batch (25--83\,\% ping loss from kratos) and its power
saving was turned off part-way, so link quality changed over time; byte counts did not. (2) Delta's
standard deviation exceeds its mean ($5.2 \pm 10.9$\,s) because run 4 took 21.7 and 27.8\,s on the two
receivers while every other run took under 1\,s; the CSVs do not show the cause. (3) Plain mode's
setup time ($0.65 \pm 1.08$\,s) comes from one 5.075\,s value (run 5, minipc); the other nine are
0.045--0.364\,s. (4) TLS received more chunk bytes than the file size: besides the excluded run 4,
defalt in run 5 received 107\,665\,543 bytes, a duplicate fetch (\#60) that \code{analyze.py} does not
flag. (5) TLS was not slower than plain; on this network the difference is inside the spread. (6) The 4.1 times loss slowdown of this batch fell to 2.1 times when rerun on the FTP night (Section~\ref{sec:ftp}): the Wi-Fi, not the protocol.
% data/README.md (lossy defalt, power saving); data/full_delta/run4.csv (delta 21.706522,
% 27.767192 s); data/tls/run5.csv (plain minipc setup_s 5.075; tls defalt chunk_bytes 107665543).

\textbf{Limitations.} All runs used one Wi-Fi network whose quality changed during the batch, so
time differences smaller than the spread mean ``no clear effect''. n = 5 runs gives wide intervals,
and two receivers are averaged per run. Times are measured on each receiver from the first
NOTIFY with a 0.5\,s polling step, so sub-second values are coarse. The chunk-size condition and
Wi-Fi vs wired were not run. Host clocks differ (kratos logs are 5\,h ahead), which is why we never
compare times across hosts. The captures in Section~\ref{sec:capture} keep only 256 bytes per
frame, so they give counts and headers, not payload timings.
% data/README.md; data/<experiment>/run*.csv (e2e_from first_notify, poll_s 0.5);
% data/captures/README.md (snaplen 256, kratos clock).

\subsection{Comparison with FTP}
\label{sec:ftp}

To test our claims against an existing protocol (PROJECT\_CONTEXT.md §5) we measured plain FTP on
the same three hosts, file and edits (21 November 2026): kratos ran vsftpd, and minipc and defalt
downloaded at the same time with Python's \code{ftplib}, checking each file's SHA-256 afterwards.
Each run has four variants: the full 100\,MiB file, the file after the 1\,KiB overwrite, after the
1\,KiB insert at offset 0, and the full file under 5\,\% netem loss on kratos. To compare times
fairly, meshsync's link condition (clean, 5\,\% loss, +50\,ms) was rerun right after the FTP batch
on the same night. Table~\ref{tab:ftp} gives mean $\pm$ sd over n = 5 runs, the two receivers
averaged per run, as computed by \code{analyze.py}.
% data/ftp/README.md (setup, variants, same-night batch, times 22:16-23:32 and 23:59-01:23);
% analysis/analyze.py on data/ftp/run*.csv and data/ftp/meshsync_same_night/run*.csv.

\begin{reptable}{>{\raggedright\arraybackslash}p{\dimexpr0.20\linewidth-2\tabcolsep\relax}>{\raggedright\arraybackslash}p{\dimexpr0.20\linewidth-2\tabcolsep\relax}>{\raggedright\arraybackslash}p{\dimexpr0.16\linewidth-2\tabcolsep\relax}>{\raggedright\arraybackslash}p{\dimexpr0.20\linewidth-2\tabcolsep\relax}>{\raggedright\arraybackslash}p{\dimexpr0.24\linewidth-2\tabcolsep\relax}}
\caption{FTP vs meshsync on the same hosts (mean $\pm$ sd, n = 5).}\label{tab:ftp}\\
\toprule
\colhead{Case} & \colhead{FTP time (s)} & \colhead{FTP bytes} & \colhead{meshsync time (s)} & \colhead{meshsync chunk bytes} \\
\midrule
\endfirsthead
\multicolumn{5}{@{}l}{\tablecontinued}\\
\toprule
\colhead{Case} & \colhead{FTP time (s)} & \colhead{FTP bytes} & \colhead{meshsync time (s)} & \colhead{meshsync chunk bytes} \\
\midrule
\endhead
\bottomrule
\endlastfoot
Full file, clean & 106.5 $\pm$ 32.5 & 104\,857\,600 & 162.8 $\pm$ 14.0$^{a}$ & 104\,874\,436 \\
Full file, 5\,\% loss & 248.9 $\pm$ 56.2 & 104\,857\,600 & 346.3 $\pm$ 59.0$^{a}$ & 104\,881\,052 \\
Loss $\div$ clean & 2.3$\times$ & & 2.1$\times$$^{a}$ & \\
1\,KiB overwrite & 88.5 $\pm$ 17.3 & 104\,857\,600 & 5.2 $\pm$ 10.9$^{b}$ & 75\,667 \\
1\,KiB insert at 0 & 87.8 $\pm$ 33.4 & 104\,858\,624 & 0.30 $\pm$ 0.41$^{b}$ & 66\,594 \\
\end{reptable}
{\footnotesize $^{a}$\,Same night as FTP. $^{b}$\,Main batch (Table~\ref{tab:experiments}): delta sync and CDC; not rerun on the FTP night, so compare bytes, not times.\par}
% FTP: analysis of data/ftp/run*.csv (full 106.5 +- 32.53, loss 248.9 +- 56.16, overwrite
% 88.5 +- 17.28, insert 87.83 +- 33.42; bytes sd 0). meshsync same night:
% data/ftp/meshsync_same_night (clean 162.8 +- 14.03, loss 346.3 +- 59.04, loss chunk bytes
% 104,881,052 +- 14,794). 248.9/106.5 = 2.34; 346.3/162.8 = 2.13.

\textbf{What the comparison shows.} After a small edit FTP downloads the whole file again: 1\,386
times the bytes meshsync sent after the overwrite, and 1\,575 times after the insert, because FTP has
no delta transfer. For a fresh full copy FTP was faster on the same night (106.5 vs 162.8\,s): one
plain TCP stream with no chunking, hashing or per-chunk requests, so less $d_\mathrm{proc}$ and fewer
round trips. meshsync's advantage is after edits, in authentication and in integrity, not in raw bulk
speed. Loss hurt both about equally on the same night (2.1 vs 2.3 times the clean time), so the 4.1
times of the main batch (Table~\ref{tab:experiments}) came from that night's worse Wi-Fi, not from
the protocol. A capture of one FTP download (not committed, because it holds the test password)
shows the weaknesses §5 lists: \code{USER} and \code{PASS} in plain text (frames 6 and 9), and two TCP
connections per transfer, the control connection on port 21 (28 frames) and a separate passive data
connection on port 40029 (77\,918 frames). meshsync sends only nonces and MACs in its handshake
(Section~\ref{sec:failure}), uses one connection per peer pair, and checks SHA-256 per chunk, while
the FTP clients had to compare SHA-256 themselves.
% 104,857,600 / 75,667 = 1,385.8; 104,858,624 / 66,594 = 1,574.6. Capture: data/ftp/README.md
% ("What the capture shows").
